\begin{table}[t]
\centering
\footnotesize
\setlength{\tabcolsep}{2.5pt}
\caption{Split-level paired gap summary for the strict-filter and no-filter heuristics against full display. The inferential unit is the split-level paired mean net-profit gap; RC uses six split pairs, while Austin and Seattle each use two split pairs and are therefore descriptive only. Austin and Seattle each have only two split pairs; their gaps are reported as descriptive checks, not stable interval estimates.}
\label{tab:policy_gap_uncertainty_summary}
\resizebox{\linewidth}{!}{%
\begin{tabular}{@{}p{0.16\linewidth}rrrrrrr@{}}
\hline
Study & Policy & Mean gap & Split-level range & Range type & Positive-split share & Gap as \% of |full| & Sampling unit \\
\hline
RC outside-option & strict-filter heuristic & 3.97 & [-6.46, 12.07] & Split-level range (6 pairs) & 0.83 & 0.09 & 6 split pairs \\
RC outside-option & no-filter heuristic & 6.61 & [1.41, 11.01] & Split-level range (6 pairs) & 0.83 & 0.15 & 6 split pairs \\
Austin & strict-filter heuristic & -6.20 & [-26.25, 13.86] & Two-pair descriptive range & 0.50 & -0.21 & 2 split pairs (descriptive) \\
Austin & no-filter heuristic & 32.12 & [24.05, 40.20] & Two-pair descriptive range & 1.00 & 1.07 & 2 split pairs (descriptive) \\
Seattle & strict-filter heuristic & 46.92 & [43.06, 50.78] & Two-pair descriptive range & 1.00 & 1.49 & 2 split pairs (descriptive) \\
Seattle & no-filter heuristic & 51.96 & [50.26, 53.65] & Two-pair descriptive range & 1.00 & 1.65 & 2 split pairs (descriptive) \\
\hline
\end{tabular}
}
\end{table}
